\documentclass[11pt,a4paper]{article}
\usepackage[T1]{fontenc}
\usepackage{lmodern}
\usepackage[margin=26mm]{geometry}
\usepackage{amsmath,amssymb,amsthm,booktabs,array,longtable}
\usepackage{microtype,multicol}
\usepackage[hidelinks]{hyperref}
\setlength{\parskip}{3pt}
\setlength{\emergencystretch}{2em}
\raggedbottom
\allowdisplaybreaks[1]
\clubpenalty=10000
\widowpenalty=10000
\newtheorem{theorem}{Theorem}[section]
\newtheorem{lemma}[theorem]{Lemma}
\newtheorem{corollary}[theorem]{Corollary}
\theoremstyle{remark}
\newtheorem{remark}[theorem]{Remark}
\newcommand{\Z}{\mathbb Z}
\newcommand{\Q}{\mathbb Q}
\newcommand{\N}{\mathbb N}
\newcommand{\F}{\mathbb F}
\newcommand{\ps}[1]{#1[[x]]}
\newcommand{\Rev}{\operatorname{Rev}}
\newcommand{\rev}{^{\langle-1\rangle}}
\newcommand{\Plus}{\operatorname{Plus}}
\newcommand{\Minus}{\operatorname{Minus}}
\newcommand{\coeff}[2]{[x^{#1}]#2}
\newcommand{\E}{\mathcal E}
\Urlmuskip=0mu plus 1mu
\hypersetup{pdftitle={A380060: Square support in an implicit rational theta equation},pdfauthor={Dev Goyal},pdfsubject={Separate manuscript with AI disclosure and Lean proof map}}
\title{A380060\\[4pt]\large Square support in an implicit rational theta equation}
\author{Dev Goyal}
\date{24 October 2026}
\begin{document}
\maketitle
\begin{center}
\fbox{\begin{minipage}{0.9\linewidth}\small
\textbf{AI disclosure.} ChatGPT (OpenAI) generated the mathematical arguments,
Lean proofs, and this manuscript in response to Dev Goyal's requests.
Dev Goyal is the sole human author. The source sequence and conjecture or
observation are credited to Paul D. Hanna.
\textbf{Independent human review is pending.}
\end{minipage}}
\end{center}
\begin{abstract}
We construct the unique integral solution of the positive-index rational theta equation $A=1+2\Sigma(A)$ and prove $A\equiv1+2\sum_{j\ge1}x^{j^2}\pmod4$. Its coefficient at a positive index is two modulo four exactly at squares and is divisible by four otherwise.
\end{abstract}

\section{Source, notation, and scope}

The source is Paul D. Hanna's OEIS entry A380060 \cite{oeis060},
retrieved on 24 October 2026. Its initial sequence index is 0.
The constant coefficient is one; the residue classification below
is restricted to positive indices.

All series are formal. The set $\N$ includes zero, and $[x^n]U$ denotes a
coefficient. A normalized series is $x+O(x^2)$; $O(x^r)$ means divisibility
by $x^r$, over any coefficient ring used below. We write $U\circ V=U(V(x))$
only when $V(0)=0$. Compositional inversion, denoted $\Rev(U)$ or $U\rev$,
is distinct from multiplicative inversion of a unit. Congruences are
coefficientwise; ordinary products and compositional iterates are distinct.

\section{Units, coefficient maps, and doubling}

A constant-one series $U=1+\sum_{n>0}u_nx^n$ has a unique multiplicative
inverse $V=1+\sum_{n>0}v_nx^n$: the coefficient equation for $UV=1$ is
\[
 v_n=-\sum_{k=1}^n u_kv_{n-k}\qquad(n>0).
\]
Thus inversion uses only ring operations, is defined integrally, and commutes
with coefficient homomorphisms. If two constant-one series agree modulo
$x^r$, the recurrence shows that their inverses agree modulo $x^r$ as well.
Locally finite coefficient sums, products, and integral powers also commute
with coefficient maps.

\begin{lemma}[Doubling]\label{lem:binary}
For formal series $U,V$ over $\Z/4\Z$, $2U=2V$ if and only if their
reductions modulo two agree.
\end{lemma}
\begin{proof}
For integers $u,v$, $4\mid2(u-v)$ is equivalent to $2\mid u-v$. Apply
this equivalence to every coefficient. No cancellation by two inside
$\Z/4\Z$ is used.
\end{proof}

\section{The literal positive-index source}

For a formal series $U$ over a commutative ring, put
\[
 Q_j(U)=\frac{U-x^j}{1-x^jU}\quad(j\ge1),\qquad
 \Sigma(U)=\sum_{j\ge1}x^{j^2}Q_j(U)^j. \tag{1}
\]
The denominator is a constant-one unit, so the quotient denotes a genuine
multiplicative inverse. The $j$th summand has order at least $j^2$.
Each coefficient is therefore a finite sum, and no analytic convergence
assumption is required. This is the positive-index formula (2) given by OEIS.
A separate formal Laurent-series equivalence with the bilateral title is
outside the claimed Lean coverage.

The selected implicit series is the solution of
\[
 A=1+2\Sigma(A). \tag{2}
\]
\begin{lemma}[Existence for the implicit theta equation]\label{lem:thetaexist}
Equation (2) has a unique formal solution over any commutative ring, and
its constant coefficient is one. Its construction commutes with
coefficient ring homomorphisms.
\end{lemma}

\begin{proof}
Suppose $U-V=O(x^r)$. For constant-one units $D,E$, the identity
\[
 D^{-1}-E^{-1}=D^{-1}(E-D)E^{-1}
\]
shows that inversion preserves coefficient agreement. Hence
$Q_j(U)-Q_j(V)=O(x^r)$, and the same holds for the difference of their
$j$th powers. Multiplication by $x^{j^2}$ improves this to at least
$O(x^{r+1})$. Since each coefficient sum is finite,
\[
 (1+2\Sigma(U))-(1+2\Sigma(V))=O(x^{r+1}).
\]
The right-hand side of (2) always has constant coefficient one. Set its constant coefficient to one. Once the coefficients below
$n$ have been chosen, the improvement of agreement shows that the degree-$n$
coefficient on the right depends only on those earlier choices. Choose the
next coefficient to equal it. Induction gives a fixed point over $K$, integral when $K=\Z$, and
the same coefficient comparison proves uniqueness. Mapping the locally finite
sums, unit inverses, and the equation gives transport by uniqueness.
\end{proof}

\begin{theorem}[A380060 modulo four]\label{thm:theta060}
The integral solution of (2) satisfies
\[
 A\equiv1+2\sum_{j\ge1}x^{j^2}\pmod4.
\]
Thus for every positive index $n$, its coefficient is two modulo four
exactly when $n$ is a square, and is divisible by four otherwise.
\end{theorem}

\begin{proof}
Equation (2) first gives $A\equiv1\pmod2$. Reduction of (1) modulo two
then gives
\[
 Q_j(A)\equiv\frac{1-x^j}{1-x^j}=1\pmod2,
 \qquad
 \Sigma(A)\equiv\sum_{j\ge1}x^{j^2}\pmod2.
\]
Doubling the latter equality lifts it to an equality modulo four by
Lemma~\ref{lem:binary}. Inserting this equality into (2) proves the
whole-series congruence. Distinct positive $j$ give distinct squares,
so its coefficient support is exactly as stated.
\end{proof}


\section{Separate Lean proof and checked scope}

The entry file is \nolinkurl{lean/Entries/A380060.lean}. Its explicit
theorem statements and proof terms use the namespace
\nolinkurl{FinalCollection.A380060}. The complete imported implementations
and their shared mathematical dependencies are included in the collection.
The following local theorem names give the proof map:

\begin{multicols}{2}\raggedright\begin{itemize}\setlength{\itemsep}{2pt}\setlength{\parsep}{0pt}
\item \nolinkurl{implicit_source_exists_unique}.
\item \nolinkurl{implicit_actual_source}.
\item \nolinkurl{a380060_statement}.
\item \nolinkurl{complete_series_mod_four}.
\end{itemize}\end{multicols}

The main implementation module is \nolinkurl{Round4.ThetaRational}.
The catalogue records the full original endpoint names and the new entry
theorems. The project pins Lean 4.19.0 and Mathlib revision
\nolinkurl{c44e0c8ee63ca166450922a373c7409c5d26b00b}. The collection
report records source recompilation and transitive axiom checks. Accepted
axioms are propositional extensionality, classical choice, and quotient
soundness. All index conclusions above are unbounded; exact finite
coefficient checks in the package provide separate corroboration.

\paragraph{Limits and status.}
The formal source is the explicit positive-index formula (2) in OEIS. A separate formal Laurent-series equivalence to the bilateral title is not claimed. No convergence or real special-value assertion is proved.
The manuscript consolidates completed results; it does not assert global
novelty, exclusive priority, peer review, OEIS acceptance, or a new submission.
Standard algebraic ingredients are not claimed as new discoveries.

\begin{thebibliography}{9}
\raggedright
\bibitem{oeis060}
Paul D. Hanna, \emph{OEIS A380060}, 6 February 2025.
\url{https://oeis.org/A380060}. Retrieved 24 October 2026.
\end{thebibliography}
\end{document}
